\section{引言}
公交运营人员分析线路运行时，需要知道耗时发生在哪些站间区间。整趟行程的总时长可以用于比较线路效率，却不能说明某次延误主要发生在前段还是后段。分段运行时间能够支持瓶颈区间筛查和分时段运行比较。然而，完整局部记录依赖持续定位与可靠的站点匹配，设备采样间隔和记录缺失会使部分区间不可观测。利用已知总时长和有限局部观测估计各段时间，因而是旅行时间推断中的实际需求\cite{hellinga2008,hofleitner2011,yin2015entryexit}。

本文研究已经完成行程的分段运行时间估计。这里的分段是有向站间区间，总时长是各区间运行时间之和，不包含停站时间。它与起终点时间戳之差有所区别：后者还可能包含停站和其他非运行时间，需要先分离这些部分并对齐区间边界。实验对公开数据的运行时间求和构造已知总时长，以考察精确总时长条件下的分段估计。

考虑一趟经过三个连续区间、运行总时长为十分钟的行程。$(2,3,5)$和$(5,3,2)$分钟是两组可行分配。训练中即使提供中间区间三分钟的标签，也无法区分这两种分配；推理时不提供当前行程的局部标签。这两组结果分别将末段和首段识别为最耗时的位置，可能导致不同的运营判断。其他行程中相同区间的记录能够帮助区分这两种分配；但如果当前行程的出发时段或路线条件发生变化，直接沿用历史平均值又可能掩盖差异。这个例子说明，模型既需要利用同一区间的重复观测，也需要判断当前行程应如何偏离平均规律。

已有研究通过自由流信息、局部时间分布或统计交通流模型为总时长分解提供依据\cite{hellinga2008,hofleitner2011}。路径不完整时，路段推断还需联合估计经过的路径与局部时间\cite{yin2015entryexit,sun2022prltte}。本文关注路径顺序已知、区间重复出现且部分单区间标签可见的条件。这里不需要重新恢复路径，关键在于如何把跨行程可共享的信息与当前行程的差异结合起来，并保证分段输出符合已知总时长。

同段均值利用重复观测，却把随出发时段和行程条件变化的耗时视为固定值；本文的直接预测基线也共享区间嵌入并使用当前行程上下文，但直接输出正值分数，不设置单独的重新分配分支。本文将两者分开建模：共享分支根据上下文估计基础时间，序列分支提出有界调整，随后按已知总时长比例缩放。本文的贡献是这一分配参数化及其在稀疏标签下的分析；其精度能否优于直接预测仍需由实验检验。

实验结果进一步限定了这一设计的作用。标签极少时，同段均值仍是有竞争力的基线；在较高标签预算下，本文方法改善了同段均值的误差，但对同输入直接预测没有稳定优势。关闭行程修正会使同一模型的误差上升，说明修正分支有助于基础估计，却不能据此推断完整方法优于直接预测。相同标签数量若只覆盖较少区间身份，误差会显著增加；在身份集合和逐身份标签数相同的控制中随机改变标签位置，也没有恢复较广覆盖时的精度。由此可见，已知总时长排除了不满足总量的分配，却不能单独判定哪一种局部时间分配正确；方法能否作出有用判断，还取决于共享区间信息的覆盖程度。
